% Generated draft: requires mathematical and source review.
\documentclass[11pt]{article}
\usepackage[margin=25mm]{geometry}
\usepackage{fontspec}
\setmainfont{Latin Modern Roman}
\usepackage{amsmath,amssymb,mathtools,unicode-math}
\setmathfont{Latin Modern Math}
\usepackage{enumitem,needspace,xurl,hyperref}
\setlength{\parindent}{0pt}
\setlength{\parskip}{5pt}
\setlength{\emergencystretch}{3em}
\newcommand{\N}{\mathbb{N}}
\newcommand{\Z}{\mathbb{Z}}
\newcommand{\Q}{\mathbb{Q}}
\newcommand{\R}{\mathbb{R}}
\newcommand{\C}{\mathbb{C}}
\newcommand{\F}{\mathbb{F}}
\newcommand{\A}{\mathbb{A}}
\newcommand{\PP}{\mathbb{P}}
\newcommand{\E}{\mathbb{E}}
\newcommand{\eps}{\varepsilon}
\newcommand{\dd}{\,\mathrm{d}}
\begin{document}
\section*{UnsolvedMath ID: 30003533}
\textbf{Source problem code:} catalogue entry
\subsection{Definitions and mathematical statement}
The supplied record asks about high-frequency coercivity of a combined boundary-integral operator \(A_k\) on strongly nontrapping, nonconvex scattering domains, mentioning smooth star-shaped domains as examples. It does not define the operator, its function space, the scattering boundary-value problem, or the intended meaning of coercivity. Consequently, a unique precise mathematical proposition cannot be recovered from this record alone.

The following is a conditional formulation template. Let \(\mathcal D\) denote the intended class of strongly nontrapping, nonconvex domains; its ambient dimension, boundary regularity and precise nontrapping condition remain unspecified. For each \(D\in\mathcal D\), suppose that a complex Hilbert space \(H_D\), its inner product, and operators \(A_k:H_D\to H_D\) have been specified for real wavenumbers \(k>0\). If coercivity is intended in the absolute numerical-range sense, define
\[
 c_D(k):=\inf_{v\in H_D\setminus\{0\}}
 \frac{\left|\langle A_kv,v\rangle_{H_D}\right|}{\|v\|_{H_D}^{2}}.
\]
Under that convention, the question has the form: for every \(D\in\mathcal D\), do there exist \(k_0(D)>0\) and a function \(\gamma_D:[k_0(D),\infty)\to(0,\infty)\) such that
\[
 \left|\langle A_kv,v\rangle_{H_D}\right|
 \geq \gamma_D(k)\|v\|_{H_D}^{2}
 \qquad
 \text{for every }k\geq k_0(D)\text{ and every }v\in H_D,
\]
with \(\gamma_D(k)\) satisfying the intended restriction of only mild degeneration as \(k\to\infty\)? The record supplies no quantitative restriction defining mild degeneration, so positivity alone does not capture that qualification.

The displayed formulation is conditional, not an identification of the intended operator or coercivity convention. A real-part convention would instead require
\[
 \operatorname{Re}\langle A_kv,v\rangle_{H_D}
 \geq \gamma_D(k)\|v\|_{H_D}^{2},
\]
possibly after a specified phase rotation; a boundary trace-space formulation could require a specified duality pairing instead of a Hilbert inner product. The record does not select among these possibilities or specify whether constants must be uniform over domains.

Smooth star-shaped domains are supplied as motivating examples, but no definition or hypothesis establishes their membership in the intended strongly nontrapping class. The supplied literature assessment states that uniform coercivity is known for convex and star-shaped formulations while the general strongly nontrapping, nonconvex question remains open. This qualification must be retained: coercivity of a particular star-shaped formulation cannot be identified with coercivity of the unspecified \(A_k\) without further information.
\subsection{Short English statement}
The question asks whether the combined boundary-integral operator remains coercive at high frequency for strongly nontrapping, nonconvex scattering domains, allowing a coercivity bound that weakens only mildly with frequency. A precise statement requires the operator, boundary-value problem, function space, coercivity convention, geometric hypotheses and permitted rate of degeneration. The supplied open-status label and claims about known convex and star-shaped formulations are dated curation, not independently verified conclusions.
\subsection{Sources}
\begin{itemize}
\item Supplied local UnsolvedMath registry record, permanent ID \(30003533\), source problem code \texttt{catalogue entry}, entitled ``High-Frequency Coercivity on Nonconvex Nontrapping Domains,'' in the Oberwolfach Reports Open Problems collection. No individual proposer is supplied; the proposed year is recorded as 2017.
\item Supplied report attribution: \emph{Material Theories (2018)}, Oberwolfach Reports, pp.~2047--2099, DOI \texttt{10.4171/owr/2017/33}. \url{https://doi.org/10.4171/owr/2017/33}. The literature-source list also labels this reference ``Oberwolfach Report 33/2017.'' The supplied 2018 title date and 2017 report designation require reconciliation; no source visit or bibliographic verification is claimed.
\item Supplied literature reference: Spence--Kamotski--Smyshlyaev, ``Combined boundary integral equations.'' \url{https://doi.org/10.1002/cpa.21543}. The record does not provide a specific theorem or operator definition from this reference.
\item Literature checked at: \texttt{2026-08-21}. This is supplied dated curation, not independently verified. Its assessment is that uniform coercivity is known for convex and star-shaped formulations, whereas the stated high-frequency question for general strongly nontrapping, nonconvex domains remains open. The background explicitly requests independent verification of time-sensitive status and cited claims.
\item The registry labels the problem \texttt{open} and its statement \texttt{corrected\_verified}; the accompanying verification describes preservation of mathematical content during repairs to notation, typography and sentence structure. These supplied labels do not establish a complete mathematical formulation or an independently verified literature status.
\end{itemize}
\end{document}
